\section{Modelling the operator's response}\label{sec:operator}

Section~\ref{sec:turn-chain} establishes partial traceability. It does not
establish what the operator attended to, why a proposal went unanswered or
whether the proposal was useful. These distinctions matter when the harness
uses follow-through to adjust future proposals.

In the audited outer-loop record, one action class received 108 unanswered
proposals, as recorded in \texttt{WR-0} (\texttt{futon3/library/war-room/}
\texttt{wr-0-organise-without-apparatus.flexiarg}). Its learned credit fell.
The record showed emissions
without a judgement, so the decrease could not distinguish rejection from
absence. The notification interface offered \texttt{nag}, \texttt{brief} and
\texttt{silent}; these describe the system's communication, not the operator's
alternative activities. A better notification policy alone would not identify
the missing response.

The proposed remedy is to record considered-and-accepted,
considered-and-declined and no-response separately, with an explicit basis for
any claim that a proposal was considered. Traces may help decide when to ask,
but cannot supply that judgement. Whether such a model improves selection
remains an intervention to test, not a result of finding the traces.

\subsection*{Reading the mechanism maps}

The following drawings answer different questions. The war-room overlay relates
constraints to mechanisms in its recorded census. The registry-generated map
shows wiring claims and recorded routes; the equation graph shows dependencies
between symbols. A dependency edge is not evidence that a call occurred.
Historical marks must be read at their source dates; current repair status is
on the issue board.

\begin{figure}[h]
\centering
\includegraphics[width=\linewidth]{aif-control-map-futon}
\caption{War-room constraints over the control map. Rings report the census in \texttt{empirics-futon/wr-overlay.edn}, including its recorded July cadence failures. They identify interfaces to investigate, not a fresh census of live execution.}\label{fig:wr-overlay}
\end{figure}
\begingroup
\addtocounter{figure}{-1}
\renewcommand{\thefigure}{\arabic{figure}A}
\begin{figure}[h]
\centering
\includegraphics[width=\linewidth]{aif-control-map-live}
\caption{Registry-generated wiring and recorded routes. The legend distinguishes observed, conditional, undrawn and retired relations; the footer pins the source. A route is an instrumented call sequence, not the dependency graph in Figure~\ref{fig:aif-dag}.}\label{fig:wr-overlay-live}
\end{figure}
\begin{figure}[h]
\centering
\includegraphics[width=\linewidth]{aif-equation-dag}
\caption{Equation dependencies. An edge identifies a symbol imported by its target; blue edges also appear in the wiring map, amber dashed edges do not. Grey nodes define no equation. This compares representations, without declaring either a witnessed execution.}\label{fig:aif-dag}
\end{figure}
\input{sec-aif-conformance-generated}
\endgroup

\subsection*{From traceability to a tested response model}

Clock-in already supplies some turn-to-mission links. The unresolved questions
are its coverage and interpretation, not the absence of every such link.
Retrieval adds candidate associations; mission prose supplies authored
citations. None of these records proves attention or pattern use. An explicit
application receipt would need to identify the actor, pattern, action and
witness, while retaining unknown or disputed attribution.

A useful experiment would compare a declared response model with the current
interpretation on the same recorded proposals, then test the effect of any
changed selection rule. It would preserve unanswered cases as such, record the
cost of eliciting a response, and distinguish a confirmed decline from an
unseen proposal. It must not infer success from a higher response rate alone.

The historical backlog map illustrates why a proposed repair also needs a
named implementation. Overlaying a mechanism gap on a mission deposit exposed
places where no mission supplied the missing behaviour. That is a comparison
of scopes, not evidence that a planned mission later ran.
\begin{figure}[h]
\centering
\includegraphics[width=\linewidth]{aif-cascade-map}
\caption{Historical mission-deposit overlay, from \texttt{empirics-futon/cascade-map.edn}. Blue rings connect proposed work to mechanism gaps. The recorded live mark is a bounded probe; it does not establish the other missions or present-day operation.}\label{fig:cascade}
\end{figure}
\begin{figure}[h]
\centering
\includegraphics[width=6.2cm]{register-key}
\caption{Evidence strength and age are separate dimensions. A record of a plan, build, run or live probe must retain when and how it was obtained. The issue board uses its own workflow columns and preserves these evidential distinctions.}\label{fig:register-key}
\end{figure}

\subsection*{What would consume the traces?}\label{tab:consumers}

Three candidate consumers need different evidence: retrieval can use textual
similarity; correction detection needs a relation between a statement and its
revision; attribution of work needs an actor's confirmation or an equivalent
witness. Joining their records is useful only if those differences survive the
join. A retrieved neighbour cannot silently become an applied pattern, and an
agent's commit cannot silently become an observation of operator attention.

The per-node dossiers that follow expose the corresponding implementation and
receipt basis. They are a way to locate and test these obligations, not a
claim that the operator-feedback loop has been closed.
